\section{Вычислительный эксперимент в математическом моделировании}
\label{sec:computational_experiment}

\subsection{Понятие и цели вычислительного эксперимента}


\textbf{Вычислительный эксперимент} --- это исследование математической
модели с помощью вычислительной техники и численных методов вместо
непосредственного проведения физического эксперимента над реальным объектом.

Общая схема математического моделирования:

\[
\text{реальный объект}
\rightarrow
\text{математическая модель}
\rightarrow
\text{численный метод}
\rightarrow
\text{компьютерная программа}
\rightarrow \\
\text{вычислительный эксперимент}
\rightarrow
\text{анализ результатов}
\rightarrow
\text{выводы}.
\]

\subsection{Отличия вычислительного эксперимента от реального}

\textbf{Объект исследования}

В реальном эксперименте исследуется непосредственно физический объект
или процесс.

В вычислительном эксперименте непосредственно исследуется
\textbf{математическая модель} объекта.

\textbf{Стоимость и безопасность}

Реальные эксперименты могут требовать: дорогостоящего оборудования; большого количества материалов; значительного времени.

Вычислительный эксперимент часто значительно дешевле и безопаснее.

Особенно это важно при исследовании опасных или разрушительных
процессов, например, исследовать
поведение конструкции при экстремальной нагрузке.

\textbf{Скорость проведения}

Компьютерный эксперимент позволяет быстро провести большое количество
расчётов.

Можно исследовать математическую модель при различных: температурах; давлениях; скоростях; материалах; размерах; начальных условиях.

Параметры можно менять автоматически и получать целую серию результатов.

\textbf{Возможность исследовать недоступные ситуации}

Например, можно моделировать: климатические процессы; движение космических аппаратов; распространение эпидемий; поведение конструкций при аварийных нагрузках; различные природные процессы.

\textbf{Воспроизводимость}

При одинаковых: исходных данных; математической модели; численном методе; параметрах алгоритма компьютерный эксперимент обычно можно повторить практически идентичным образом.

В реальном эксперименте всегда присутствуют случайные факторы,
погрешности измерительных приборов и внешние воздействия.


\textbf{Погрешности}

Важно понимать, что компьютерный эксперимент не гарантирует
абсолютно правильный результат.

Основными источниками погрешности являются: упрощение математической модели;  неточность исходных данных; приближённость численного метода; слишком крупный шаг сетки; ошибки в программе; точность вычислений.

Следовательно, можно получить очень точное решение
\textbf{неправильной математической модели}.

Поэтому высокая точность вычислений сама по себе ещё не означает
достоверность результата.

\subsection{Этапы вычислительного эксперимента}

\textbf{Постановка физической задачи}

Пусть имеется стержень длиной $L$.

Температура внутри стержня распределена неоднородно. Необходимо
определить температуру $T(x,t)$ в каждой точке стержня в каждый
момент времени.

\textbf{Построение математической модели}

Процесс можно описать одномерным уравнением теплопроводности:

\[
\frac{\partial T}{\partial t}
=
\alpha
\frac{\partial^2 T}{\partial x^2},
\]

где:

\begin{itemize}
	\item $T(x,t)$ --- температура;
	\item $x$ --- координата вдоль стержня;
	\item $t$ --- время;
	\item $\alpha$ --- коэффициент температуропроводности материала.
\end{itemize}

Кроме уравнения необходимо задать начальные и граничные условия.

Начальное условие:

\[
T(x,0)=T_0.
\]

Это означает, что в начальный момент весь стержень имеет температуру
$T_0$.

На левом конце поддерживается температура $T_1$:

\[
T(0,t)=T_1.
\]

На правом конце зададим условие теплоизоляции:

\[
\frac{\partial T}{\partial x}(L,t)=0.
\]

Таким образом, физическая задача преобразуется в математическую.

\textbf{Численное решение}

В общем случае решить полученное уравнение аналитически может быть
сложно или невозможно. Поэтому используется численный метод.

Разобьём стержень на сетку:

\[
x_0,x_1,\ldots,x_N
\]

с шагом

\[
\Delta x.
\]

Время также разобьём на шаги:

\[
t_0,t_1,\ldots,t_M
\]

с шагом

\[
\Delta t.
\]

Вместо непрерывной функции $T(x,t)$ компьютер будет рассчитывать
приближённые значения

\[
T_i^n \approx T(x_i,t_n).
\]

Например, вторую производную можно заменить конечной разностью:

\[
\frac{\partial^2 T}{\partial x^2}
\approx
\frac{
	T_{i+1}^n
	-
	2T_i^n
	+
	T_{i-1}^n
}{
	(\Delta x)^2
}.
\]

В результате исходное дифференциальное уравнение превращается
в систему вычислительных операций, которые компьютер может выполнить.

\textbf{Проведение вычислительного эксперимента}

После построения алгоритма можно провести серию расчётов.

Например, сначала зададим

\[
\alpha=\alpha_1.
\]

Получим распределение температуры через:

\[
t=10,\quad 20,\quad 30\text{ с}.
\]

Затем изменим параметр:

\[
\alpha=\alpha_2
\]

и повторим расчёт.

Также можно исследовать влияние:

\begin{itemize}
	\item начальной температуры;
	\item температуры нагретого конца;
	\item длины стержня;
	\item коэффициента теплопроводности;
	\item времени;
	\item шага сетки;
	\item временного шага.
\end{itemize}

Таким образом, вычислительный эксперимент позволяет исследовать
не только вопрос

\begin{quote}
	«Какова температура?»
\end{quote}

но и вопрос

\begin{quote}
	«Как изменится результат при изменении параметров модели?»
\end{quote}


\subsection{Что важно учитывать при проведении вычислительного эксперимента}


\textbf{Адекватность математической модели}

Модель должна достаточно хорошо описывать реальный объект или процесс.

Если модель не соответствует реальности, даже очень точные вычисления не дадут достоверного результата.

\textbf{Корректность исходных данных}

Необходимо правильно задать:

\begin{itemize}
	\item параметры объекта;
	\item начальные и граничные условия;
	\item физические коэффициенты.
\end{itemize}

Ошибки в исходных данных приводят к ошибочным результатам.

\textbf{Выбор численного метода}

Для решения задачи необходимо выбрать подходящий численный метод, например:

\begin{itemize}
	\item метод конечных разностей;
	\item метод конечных элементов;
	\item метод конечных объёмов;
	\item методы Рунге--Кутты.
\end{itemize}

Выбор метода зависит от особенностей математической задачи.

\textbf{Устойчивость}

Численный метод должен быть устойчивым. При слишком большом шаге
$\Delta x$ или $\Delta t$ решение может начать сильно отклоняться
от правильного.

\textbf{Сходимость}

При уменьшении шага сетки или временного шага численное решение
должно приближаться к точному решению.

Например:

\[
\Delta x=0.1,\qquad
\Delta x=0.05,\qquad
\Delta x=0.025.
\]

Если результаты при этом стабилизируются, можно говорить о сходимости.

\textbf{Оценка погрешности}

Необходимо учитывать погрешности:

\begin{itemize}
	\item математической модели;
	\item исходных данных;
	\item численного метода;
	\item программной реализации;
	\item компьютерных вычислений.
\end{itemize}

\textbf{Верификация и валидация}

\textbf{Верификация} --- проверка правильности реализации
математической модели, алгоритма и программы.

\textbf{Валидация} --- проверка того, насколько математическая модель
соответствует реальному объекту.

Например, результаты моделирования можно сравнить с результатами
реального физического эксперимента.

\newpage